\documentclass[11pt]{article}

\usepackage[margin=1in]{geometry}
\usepackage{amsmath, amssymb, amsthm}
\usepackage{graphicx}
\usepackage{hyperref}
\usepackage{enumitem}
\usepackage{xcolor}

\title{
    \normalsize IDs: 213851579, 213310485 \\[1em] % This puts IDs at the top
    \huge Video Processing -- HW1 \\ 
    \large Harris Corner Detection and Intro to Video Processing
}

% Leave author empty or use it for names if needed
\author{} 

\date{Course 0512.4263 \quad 2025/2026}

\begin{document}

% ========================================================================
% REQUIRED: first line of the PDF must contain the submitters' IDs.
% ========================================================================

\maketitle

% ========================================================================
\section*{Section 2 / Part 2.1 -- Harris Corner Detector: Open Questions}
% ========================================================================

\subsection*{Q1. Is the Harris corner detector invariant to translation?}
\textbf{Answer:} Yes.

\medskip
\noindent\textbf{Explanation.}
Every step of the detector (gradients, pointwise products, convolution with $g$, the response $R$, NMS, and threshold) is shift-invariant: it depends only on relative pixel positions. So shifting the image by $(dx,dy)$ shifts the corner map by the same $(dx,dy)$, and the same physical points are detected.



\subsection*{Q2. Is the Harris corner detector invariant to rotation?}
\textbf{Answer:} Yes.

\medskip
\noindent\textbf{Explanation.}
The Harris response depends only on the eigenvalues of
\[
    M = \begin{bmatrix} S_{xx} & S_{xy} \\ S_{xy} & S_{yy} \end{bmatrix},
\]
since $R = \lambda_1 \lambda_2 - k(\lambda_1 + \lambda_2)^2$.
Rotating the image transforms $M$ into $R_\theta M R_\theta^{T}$, which has the same eigenvalues. So $R$ is unchanged (just evaluated at the rotated pixel location), and the same points are detected as corners.



\subsection*{Q3. Is the Harris corner detector invariant to constant illumination
($I_{out} = a \cdot I_{in} + b$)?}
\textbf{Answer:} No.

\medskip
\noindent\textbf{Explanation.}
The bias $b$ cancels in the gradients, so it has no effect. The gain $a$ scales the gradients by $a$, hence $S_{xx}, S_{yy}, S_{xy}$ by $a^2$ and the response by $R' = a^4 R$. NMS picks the same local maxima, but the fixed threshold passes a different set of pixels, so the binary corner map changes. Invariant to $b$, not invariant to $a$ $\Rightarrow$ not fully invariant.
\newpage
% ========================================================================
\section*{Section 3 -- Q3: Sobel Operator}
% ========================================================================

\noindent\textit{Describe the Sobel operator and explain how it may be used to
obtain the edge map of each frame of the video.}

\medskip
\noindent\textbf{Definition.}
% Write the definition here. Include the two 3x3 kernels:
\[
G_x =
\begin{bmatrix}
-1 & 0 & +1 \\
-2 & 0 & +2 \\
-1 & 0 & +1
\end{bmatrix},
\qquad
G_y =
\begin{bmatrix}
-1 & -2 & -1 \\
\phantom{-}0 & \phantom{-}0 & \phantom{-}0 \\
+1 & +2 & +1
\end{bmatrix}.
\]

\noindent\textbf{Gradient magnitude / direction.}
\[
|G| = \sqrt{G_x^{\,2} + G_y^{\,2}}, \qquad
\theta = \operatorname{atan2}(G_y,\, G_x).
\]

\medskip
By applying the sobel operator on a frame of the video, we get a 2-dimensional map of the gradient at each point, compute the magnitude $|G|$ and threshold. Pixels with $|G|$ above the threshold are marked as edges; the rest are background. Repeating this on every frame produces an edge-map video.

\end{document}
